\documentclass[10pt,a4paper]{amsart}
\usepackage[margin=19mm]{geometry}
\usepackage{amsmath,amssymb,mathtools}
\usepackage[scr=rsfs,cal=euler]{mathalfa}
\usepackage{xfrac}
\usepackage[unicode,hidelinks,bookmarks=false]{hyperref}
\newcommand{\ie}{i.e.\ }
\newcommand{\cf}{cf.\ }
\newcommand{\resp}{respectively\ }
\newcommand{\Subsection}{\subsection}
\providecommand{\nobreakdash}{\nobreak\hbox{-}}
\setlength{\emergencystretch}{2em}
\multlinegap0pt
\usepackage{bm}
\usepackage{bbm}
\usepackage{dsfont}
\usepackage{xspace}
\usepackage{cancel}
\usepackage{mathtools}
\usepackage{amsthm}
\makeatletter
\@ifundefined{thm}{\newtheorem{thm}{Thm}}{}
\@ifundefined{lemma}{\newtheorem{lemma}[thm]{Lemma}}{}
\makeatother
\newcommand{\escal}[1]{\langle#1\rangle}
\title[The algebraic square of an irreducible complex spinor]{The algebraic square of an irreducible complex spinor}
\author{Alejandro Gil-Garc\'ia, C. S. Shahbazi}
\date{Source: arXiv:2510.13801v1 (CC BY 4.0)}
\begin{document}
\maketitle

\noindent\emph{Source and licence.} The algebraic square of an irreducible complex spinor. Version arXiv:2510.13801v1 (CC BY 4.0), distributed under the Creative Commons Attribution 4.0 International licence, as recorded in the arXiv licence metadata for this exact version and shown on the abstract page https://arxiv.org/abs/2510.13801v1. Full rights notice: licenses/arxiv-2510.13801-CC-BY-4.0.txt.\par\smallskip

\noindent\textit{Editorial scope.} Counted unit: the Lemma whose source label is \texttt{None} in the article (source0.tex lines 3953--3958, source label \texttt{None}), delivered together with the article's own complete original proof of it (source0.tex lines 3960--4013, one \texttt{proof} environment of 54 lines). No mathematics is written, repaired or re-derived: statement and proof are the article's own text line for line. Required context retained uncounted: eq:sesquilinear6d (lines 3916--3919). Every definition, display and label of those lines is retained unchanged. Omitted from this package and not redistributed: 1--3915, 3920--3952, 3959--3959, 4014--4533. Nothing inside a retained excerpt is omitted or altered: every retained body is delivered exactly as the article's lines are written, so no omission receipt of any kind accompanies this package. Citations: the retained excerpts carry no citation command at all, so no bibliography entry of the article is transcribed and no reference is redistributed. Reference adaptation: none was needed and none was made; every reference inside the delivered unit resolves inside it, so no unresolved reference or citation is produced. Printed slips kept literal: no printed word, symbol, hypothesis or number of the counted statement or of its proof was altered. Layout and class: the article's own class is \texttt{amsart} with packages including geometry, bm, latexsym, amsmath, amstext, amssymb, amsfonts, amscd, array, multirow, amsbsy, mathrsfs, amsthm, t1enc; the wrapper is the lane's pinned \texttt{amsart} editorial wrapper at 10pt on A4, which restates the theorem-like environments the retained text needs and adds the article's own macro definitions that the retained text needs (\texttt{\textbackslash escal}), with extra wrapper packages (bm, bbm, dsfont, xspace, cancel, mathtools). The delivered numbering is the unit's own. Assets: no retained span contains a figure environment, a graphics-inclusion command, a PDF-page inclusion, a table or a TikZ picture; no image file is copied into this package, the package declares no asset dependency and no image file is redistributed with it. Licence: version arXiv:2510.13801v1 of this article is distributed under the Creative Commons Attribution 4.0 International licence, as recorded in the arXiv licence metadata for this exact version and shown on the abstract page https://arxiv.org/abs/2510.13801v1. A full rights notice is delivered at licenses/arxiv-2510.13801-CC-BY-4.0.txt. Characters: every retained excerpt is pure ASCII, so the delivered engine, XeLaTeX with its default Latin Modern text font, reports no missing character.
\begin{eqnarray}
\label{eq:sesquilinear6d}
\tau(\widehat{\alpha})=\overline{\widehat{\alpha}}\, ,\qquad   \ast (\pi\circ\tau)(\widehat{\alpha}) =   \mu\widehat{\alpha}\, , \qquad \widehat{\alpha} \diamond \beta \diamond\widehat{\alpha}= 8 (\widehat{\alpha}\diamond\beta)^{(0)}\widehat{\alpha}
\end{eqnarray}

\begin{lemma}
Using the notation introduced above, let $\widehat{\alpha} = u + i \rho  - \mu \ast u$ be the Hermitian square of an irreducible complex and chiral spinor $\eta\in \Sigma^{\mu}\setminus \{ 0 \}$. Then:
\begin{equation*}
\escal{u,u}= 0 \, , \qquad u\wedge \rho = 0\, , \qquad \rho(u^\sharp) = 0.
\end{equation*}
\end{lemma}

\begin{proof}
Suppose $u\neq 0$ and set $\beta = u\in V^*$ in the third equation of \eqref{eq:sesquilinear6d}. We compute:
\begin{equation*}
\widehat{\alpha}\diamond u=\escal{u,u}+i(\rho\wedge u+\rho(u^\sharp))+\mu\escal{u,u}\nu.
\end{equation*}

\noindent
From this, we obtain that the third equation in \eqref{eq:sesquilinear6d} is equivalent to:
\begin{equation}
\label{eq:3rdII}
\widehat{\alpha}\diamond u\diamond\widehat{\alpha}=2\escal{u,u}(u+2i\rho-\mu*u)-\rho\diamond u\diamond\rho=8\escal{u,u}(u + i \rho  - \mu \ast u)=8 (\widehat{\alpha}\diamond u)^{(0)}\widehat{\alpha}.
\end{equation}

\noindent
Regarding the three-form $\rho$, we have two possibilities: either $\rho = 0$ or $\rho \neq 0$. If $\rho = 0$ then the previous equation reduces to $\escal{u,u} = 0$. On the other hand, if $\rho\neq  0$ then the imaginary part of Equation \eqref{eq:3rdII} gives $\escal{u,u}\rho = 0$ and hence $\escal{u,u} = 0$ again. Hence $\escal{u,u} = 0$ and plugging this condition into Equation \eqref{eq:3rdII}, we obtain:
\begin{equation*}
\rho \diamond u\diamond \rho   =  2 \rho(u^\sharp) \diamond \rho = 2 (\rho(u^\sharp) \wedge \rho - \rho(u^\sharp) \triangle_1 \rho - \rho(u^\sharp) \triangle_2 \rho) = 0
\end{equation*}
 
\noindent
or, equivalently:
\begin{equation*}
\rho(u^\sharp) \wedge \rho  = 0\, , \qquad  \rho(u^\sharp) \triangle_1 \rho = 0\, , \qquad  \rho(u^\sharp) \triangle_2 \rho = 0.
\end{equation*}

\noindent
To proceed further, using that we have proven $u$ to be isotropic,  we choose a one-form $v\in V^*$ \emph{conjugate} to $u$. That is, $v$ is a nowhere vanishing isotropic one-form satisfying $\langle u , v \rangle = 1$. Hence:
\begin{equation*}
V  = \langle \mathbb{R} u^{\sharp} \rangle \oplus \langle \mathbb{R} v^{\sharp} \rangle \oplus (\langle \mathbb{R} u^{\sharp} \rangle \oplus \langle \mathbb{R} v^{\sharp} \rangle)^{\perp},
\end{equation*}

\noindent
where $V_{uv}:=(\langle \mathbb{R} u^{\sharp} \rangle \oplus \langle \mathbb{R} v^{\sharp} \rangle)^{\perp} \subset V$ denotes the orthogonal complement of $\langle \mathbb{R} u^{\sharp} \rangle \oplus \langle \mathbb{R} v^{\sharp} \rangle \subset V$. With respect to the previous splitting, we can decompose $\rho\in\wedge^3V^*$ as follows:
\begin{equation*}
\rho=u\wedge v\wedge\phi+u\wedge\omega+v\wedge\varrho+\rho^\perp,
\end{equation*}

\noindent
where $\phi\in V_{uv}^*$, $\omega,\varrho\in\wedge^2V_{uv}^*$, and $\rho^\perp\in\wedge^3V_{uv}^*$. Note that the metric $h$ on $V$ induces a positive-definite metric $h_{uv}$ on $V_{uv}$. Moreover, we define a volume form $\nu_{uv}$ on $V_{uv}$ by the relation $\nu=u\wedge v\wedge \nu_{uv}$, and consequently a Hodge star operator $*_{uv}$ on $V_{uv}$. Therefore, the duality condition $*\rho=\mu\rho$ becomes:
\begin{equation*}
*_{uv}\omega=-\mu\omega,\qquad *_{uv}\varrho=\mu\varrho,\qquad *_{uv}\rho^\perp=\mu\phi.
\end{equation*}

\noindent
We then have $\rho(u^\sharp)=-u\wedge\phi+\varrho$ and the equation $\rho(u^\sharp)\wedge\rho=0$ becomes: $$\phi\wedge\varrho=0,\qquad \phi\wedge\rho^\perp-\omega\wedge\varrho=0,\qquad \varrho\wedge\varrho=0.$$

Using that $*_{uv}\varrho=\mu\varrho$, the third equation above implies $\escal{\varrho,\varrho}=0$, which in turn implies $\varrho=0$ since the metric $h_{uv}$ on $V_{uv}$ is positive-definite. Using $*_{uv}\rho^\perp=\mu\phi$, the second equation above implies that $\rho^\perp=0$, so $\phi=0$ and $\rho=u\wedge\omega$. This immediately gives:
\begin{equation*}
\rho(u^\sharp)=0,\qquad u\wedge\rho=0
\end{equation*}

\noindent
and thus we conclude.
\end{proof}

\end{document}
